% GENERATED FILE. Edit canonical lessons, then run npm run generate:books.
% canonical-source-hash: fnv1a64:f6d49a5ab1c3204f
% canonical-lessons: FR-C318-feu-d-artifice, FR-C318-petard, FR-C318-bal, FR-C318-guirlande, FR-C318-decorer

\chapter{Le feu d'artifice}
\label{ch:fr-le-feu-d-artifice}
\hlchaptermodality{318}

\begin{chapteropening}
\textbf{By the end of this chapter:} \emph{I can describe a night of celebration: fireworks, firecrackers, the dance, garlands and decorating.}

Complete the last lesson of chapter 318: I can describe a night of celebration: fireworks, firecrackers, the dance, garlands and decorating.
\end{chapteropening}

\section[le feu d'artifice]{le feu d'artifice — the fireworks (display)}

\label{lesson:FR-C318-feu-d-artifice}



\begin{quote}
Before the new one: say the French for the float, then the French for to dress up.
\end{quote}

\subsection*{You'll want to know: le feu d'artifice}
\textbf{le feu d'artifice} — \textquotedblleft{}the fireworks (display)\textquotedblright{}.

\begin{quote}
Le feu d'artifice a commencé à vingt-trois heures.
\end{quote}

\begin{quote}
\emph{The fireworks started at eleven at night.}
\end{quote}

Plural: \emph{les feux d'artifice}.

\begin{cousinweb}
\emph{feu} is inherited from Latin \emph{focus} \textquotedblleft{}hearth\textquotedblright{} (EtymDB), and \emph{artifice} is a learned borrowing from Latin \emph{artificium} \textquotedblleft{}craft, skill\textquotedblright{}, from \emph{ars} and \emph{facere} (Lewis \& Short).
\end{cousinweb}

\begin{grammarlens}[title={What you've built}]
One more word for telling, explaining and describing at B1.
\end{grammarlens}

\subsection*{Your turn}
\begin{itemize}
\raggedright
  \item \emph{Say it:} \emph{le feu d'artifice}
  \item \emph{Say it:} \emph{le feu d'artifice}, once more
  \item \emph{Say it:} a sentence of your own with \emph{le feu d'artifice}
  \item \emph{Recall:} say the French for the float, then the French for to dress up, then say \emph{le feu d'artifice} again
\end{itemize}

\begin{tcolorbox}[breakable,colback=teal!4,colframe=teal!35!black,title={Before you move on}]
What does le feu d'artifice mean? (\textquotedblleft{}The fireworks (display)\textquotedblright{}.) Say it once more.
\end{tcolorbox}
\section[le pétard]{le pétard — the firecracker, the banger}

\label{lesson:FR-C318-petard}



\begin{quote}
Before the new one: say the French for to dress up, then the French for the fireworks.
\end{quote}

\subsection*{You'll want to know: le pétard}
\textbf{le pétard} — \textquotedblleft{}the firecracker, the banger\textquotedblright{}.

\begin{quote}
Les enfants ont lancé des pétards dans la rue.
\end{quote}

\begin{quote}
\emph{The children set off bangers in the street.}
\end{quote}

Informal: \emph{être en pétard} — to be furious.

\begin{cousinweb}
\emph{pétard} is built on \emph{péter} \textquotedblleft{}to go bang, to break wind\textquotedblright{} with the suffix \emph{-ard} (EtymDB).
\end{cousinweb}

\begin{grammarlens}[title={What you've built}]
One more word for telling, explaining and describing at B1.
\end{grammarlens}

\subsection*{Your turn}
\begin{itemize}
\raggedright
  \item \emph{Say it:} \emph{le pétard}
  \item \emph{Say it:} \emph{le pétard}, once more
  \item \emph{Say it:} a sentence of your own with \emph{le pétard}
  \item \emph{Recall:} say the French for to dress up, then the French for the fireworks, then say \emph{le pétard} again
\end{itemize}

\begin{tcolorbox}[breakable,colback=teal!4,colframe=teal!35!black,title={Before you move on}]
What does le pétard mean? (\textquotedblleft{}The firecracker, the banger\textquotedblright{}.) Say it once more.
\end{tcolorbox}
\section[le bal]{le bal — the dance, the ball}

\label{lesson:FR-C318-bal}



\begin{quote}
Before the new one: say the French for the fireworks, then the French for the firecracker.
\end{quote}

\subsection*{You'll want to know: le bal}
\textbf{le bal} — \textquotedblleft{}the dance, the ball\textquotedblright{}.

\begin{quote}
Les jeunes dansent au bal du village.
\end{quote}

\begin{quote}
\emph{The young people dance at the village dance.}
\end{quote}

\emph{Le bal populaire} — the open-air public dance, a summer tradition.

\begin{cousinweb}
\emph{bal} is built on Old French \emph{baller} \textquotedblleft{}to dance\textquotedblright{}, which is inherited from Late Latin \emph{ballāre} (EtymDB).
\end{cousinweb}

\begin{grammarlens}[title={What you've built}]
One more word for telling, explaining and describing at B1.
\end{grammarlens}

\subsection*{Your turn}
\begin{itemize}
\raggedright
  \item \emph{Say it:} \emph{le bal}
  \item \emph{Say it:} \emph{le bal}, once more
  \item \emph{Say it:} a sentence of your own with \emph{le bal}
  \item \emph{Recall:} say the French for the fireworks, then the French for the firecracker, then say \emph{le bal} again
\end{itemize}

\begin{tcolorbox}[breakable,colback=teal!4,colframe=teal!35!black,title={Before you move on}]
What does le bal mean? (\textquotedblleft{}The dance, the ball\textquotedblright{}.) Say it once more.
\end{tcolorbox}
\section[la guirlande]{la guirlande — the garland, the string of lights}

\label{lesson:FR-C318-guirlande}



\begin{quote}
Before the new one: say the French for the firecracker, then the French for the dance.
\end{quote}

\subsection*{You'll want to know: la guirlande}
\textbf{la guirlande} — \textquotedblleft{}the garland, the string of lights\textquotedblright{}.

\begin{quote}
Il y a des guirlandes dans toutes les rues.
\end{quote}

\begin{quote}
\emph{There are strings of lights in every street.}
\end{quote}

\emph{Une guirlande électrique} — fairy lights.

\begin{cousinweb}
\emph{guirlande} was borrowed from Italian \emph{ghirlanda} \textquotedblleft{}garland\textquotedblright{} (EtymDB); the origin of the Italian word is not settled.
\end{cousinweb}

\begin{grammarlens}[title={What you've built}]
One more word for telling, explaining and describing at B1.
\end{grammarlens}

\subsection*{Your turn}
\begin{itemize}
\raggedright
  \item \emph{Say it:} \emph{la guirlande}
  \item \emph{Say it:} \emph{la guirlande}, once more
  \item \emph{Say it:} a sentence of your own with \emph{la guirlande}
  \item \emph{Recall:} say the French for the firecracker, then the French for the dance, then say \emph{la guirlande} again
\end{itemize}

\begin{tcolorbox}[breakable,colback=teal!4,colframe=teal!35!black,title={Before you move on}]
What does la guirlande mean? (\textquotedblleft{}The garland, the string of lights\textquotedblright{}.) Say it once more.
\end{tcolorbox}
\section[décorer]{décorer — to decorate}

\label{lesson:FR-C318-decorer}



\begin{quote}
Before the new one: say the French for the dance, then the French for the garland.
\end{quote}

\subsection*{You'll want to know: décorer}
\textbf{décorer} — \textquotedblleft{}to decorate\textquotedblright{}.

\begin{quote}
Les enfants ont décoré la salle avec des ballons.
\end{quote}

\begin{quote}
\emph{The children decorated the room with balloons.}
\end{quote}

Also to award a medal: \emph{il a été décoré}.

\begin{cousinweb}
\emph{décorer} is a learned borrowing from Latin \emph{decorāre} \textquotedblleft{}to adorn\textquotedblright{}, from \emph{decus} \textquotedblleft{}ornament, honour\textquotedblright{} (Lewis \& Short).
\end{cousinweb}

\begin{grammarlens}[title={What you've built}]
One more word for telling, explaining and describing at B1.
\end{grammarlens}

\subsection*{Your turn}
\begin{itemize}
\raggedright
  \item \emph{Say it:} \emph{décorer}
  \item \emph{Say it:} \emph{décorer}, once more
  \item \emph{Say it:} a sentence of your own with \emph{décorer}
  \item \emph{Recall:} say the French for the dance, then the French for the garland, then say \emph{décorer} again
\end{itemize}

\begin{tcolorbox}[breakable,colback=teal!4,colframe=teal!35!black,title={Before you move on}]
What does décorer mean? (\textquotedblleft{}To decorate\textquotedblright{}.) Say it once more.
\end{tcolorbox}
